\begin{theorem}[Theorem B]\label{thm:16.7}
	Let $m,n\in\N$ with $m\ge 2$. Let $\lambda\vdash n$ and $k\in\{0,1,\dotsc,n-1\}$. Then
	\[ \rho^\lambda_m/\chi^{(1^{n-k})} = \sum_{\beta\vdash k} \rho^\beta_m\boxtimes \rho^{\lambda'/\beta'}_{m-1}. \]
\end{theorem}

\begin{proof}
	Following the notation for $\underline{t}$ in Corollary~\ref{cor:16.2} and letting $r=l(\lambda)$, observe that
	\begin{align*}
		\sum_{\gamma\vdash n} K_{\gamma,\lambda} &\cdot \big( \rho^\gamma_m/\chi^{(1^{n-k})} \big) = \Big( \sum_{\gamma\vdash n} K_{\gamma,\lambda}\cdot \rho^\gamma_m \Big)/\chi^{(1^{n-k})} \overset{\normalfont\tiny\text{\eqref{eqn:15.2}}}{=} \Big( \operatorname*{\boxtimes}_{i=1}^r \rho^{(\lambda_i)}_m\Big)/\chi^{(1^{n-k})}\\
		&\overset{\normalfont\tiny\text{(Corollary~\ref{cor:16.2}\ref{591})}}{=} \sum_{\underline{t}} \operatorname*{\boxtimes}_{i=1}^r \big( \rho^{(\lambda_i)}_m/\chi^{(1^{\lambda_i-t_i})} \big)\\
		&\overset{\normalfont\tiny\text{Lemma~\ref{lem:16.6}}}{=} \sum_{\underline{t}} \operatorname*{\boxtimes}_{i=1}^r \big( \rho^{(t_i)}_m \boxtimes \rho^{(1^{\lambda_i-t_i})}_{m-1} \big)\\
		&\overset{\normalfont\tiny\text{\eqref{eqn:15.2}}}{=} \sum_{\underline{t}} \sum_{\beta\vdash k} c^\beta_{(t_1),\dotsc,(t_r)}\cdot\rho^\beta_m \boxtimes \big(\operatorname*{\boxtimes}_{i=1}^r \rho^{(1^{\lambda_i-t_i})}_{m-1}\big)\\
		&\overset{\normalfont\tiny \text{Lemma~\ref{lem:cX}}}{=} \sum_{\beta\vdash k} \rho^\beta_m \boxtimes \cX\Big(\triv_{S_{m-1}}; \sum_{\underline{t}} c^\beta_{(t_1),\dotsc,(t_r)}\cdot \operatorname*{\boxtimes}_{i=1}^r \chi^{(1^{\lambda_i-t_i})} \Big)\up_{S_{m-1}\wr S_{n-k}}^{S_{(m-1)(n-k)}}\\
		&\overset{\normalfont\tiny\text{Corollary~\ref{cor:16.2}\ref{592}}}{=} \sum_{\beta\vdash k} \rho^\beta_m \boxtimes \cX \Big( \triv_{S_{m-1}}; \big( \operatorname*{\boxtimes}_{i=1}^r \chi^{(1^{\lambda_i})}\big) /\chi^{\beta'} \Big) \up_{S_{m-1}\wr S_{n-k}}^{S_{(m-1)(n-k)}}\\
		&\overset{\normalfont\tiny\text{\eqref{eqn:sign}}}{=} \sum_{\beta\vdash k} \rho^\beta_m \boxtimes \cX \Big( \triv_{S_{m-1}}; \big( \zeta^\lambda\cdot\sgn_{S_n} \big) /\chi^{\beta'} \Big) \up_{S_{m-1}\wr S_{n-k}}^{S_{(m-1)(n-k)}}\\
		&\overset{\normalfont\tiny\text{\eqref{eqn:sign} and \eqref{eqn:kostka}}}{=} \sum_{\beta\vdash k} \rho^\beta_m \boxtimes \cX \Big( \triv_{S_{m-1}}; \big( \sum_{\gamma\vdash n} K_{\gamma,\lambda}\cdot \chi^{\gamma'} \big) /\chi^{\beta'} \Big) \up_{S_{m-1}\wr S_{n-k}}^{S_{(m-1)(n-k)}} \\
		&= \sum_{\gamma\vdash n} K_{\gamma,\lambda} \cdot \Big( \sum_{\beta\vdash k} \rho^\beta_m \boxtimes \rho^{\gamma'/\beta'}_{m-1} \Big)\quad\text{since $\chi^{\gamma'}/\chi^{\beta'}=\chi^{\gamma'/\beta'}$ by Notation~\ref{not:rho}\ref{561}}.
	\end{align*}
	Since the matrix $(K_{\gamma,\lambda})_{\gamma,\lambda\vdash n}$ is invertible (in fact unitriangular if the partitions are ordered lexicographically, see e.g.~\cite[Chapter 2]{JK}), we deduce that $\rho^\gamma_m/\chi^{(1^{n-k})} = \sum_{\beta\vdash k} \rho^\beta_m \boxtimes \rho^{\gamma'/\beta'}_{m-1}$ for each $\gamma\vdash n$.
\end{proof}

\begin{lemma}\label{lem:cX}
	Let $m,r,a_1,\dotsc,a_r\in\N$, and let $n=\sum_{i=1}^r a_i$. For each $i\in\{1,\dotsc,r\}$, let $\nu_i\vdash a_i$. Then $\boxtimes_{i=1}^r \rho_m^{\nu_i} = \cX(\triv_{S_m};\boxtimes_{i=1}^r \chi^{\nu_i} )\up_{S_m\wr S_n}^{S_{mn}}$.
\end{lemma}

\begin{proof}
	The case $r=1$ follows from Notation~\ref{not:rho}\ref{562}. For each of notation we prove the statement for $r=2$; the case of general $r$ follows by an analogous argument. In fact, we can prove more generally that if $a,b\in\N$ and $\phi_1\in\Ch(S_a), \phi_2\in\Ch(S_b)$, then $\cX_{left}=\cX_{right}$ where
	\[ \cX_{left}\coloneqq  \left[\cX(\triv_{S_m}; \phi_1)\up_{S_m\wr S_a}^{S_{ma}} \times \cX(\triv_{S_m}; \phi_2)\up_{S_m\wr S_b}^{S_{mb}}\right]\up_{S_{ma}\times S_{mb}}^{S_{m(a+b)}} \]
	and
	\[ \cX_{right}\coloneqq  \cX\left(\triv_{S_m}; (\phi_1\times \phi_2)\up_{S_a\times S_b}^{S_{a+b}}\right)\up_{S_m\wr S_{a+b}}^{S_{m(a+b)}}, \]
	from which we recover the case of $r=2$ by setting $\phi_i=\chi^{\nu_i}$. 
	
	To prove that $\cX_{left}=\cX_{right}$, we first observe that $S_m\wr(S_a\times S_b) = S_m\wr S_a \times S_m\wr S_b$ (viewing $S_a\times S_b$ as a subgroup of $S_{a+b}$). Calling this group $U$, it is a subgroup of both $T_d\coloneqq S_{ma}\times S_{mb}$ and $T_w\coloneqq S_m\wr S_{a+b}$, and both $T_d$ and $T_w$ are subgroups of $S\coloneqq S_{m(a+b)}$. Now, by Lemma~\ref{lem:infl-ind} and the definition of $\cX(-;-)$,
	\begin{multline*}
		\cX(\triv_{S_m}; (\phi_1 \times \phi_2)\up_{S_a\times S_b}^{S_{a+b}}) = \cX(\triv_{S_m}; \phi_1\times\phi_2)\up_U^{T_w} = \left(\Infl_{S_a\times S_b}^{T_w}(\phi_1\times\phi_2)\right)\up_U^{T_w}\\
		= \left( \Infl_{S_a}^{S_m\wr S_a}(\phi_1)\times \Infl_{S_b}^{S_m\wr S_b}(\phi_2) \right)\up_U^{T_w} = \Big( \cX(\triv_{S_m};\phi_1)\times \cX(\triv_{S_m};\phi_2) \Big)\up_U^{T_w}.
	\end{multline*}
	Therefore
	\begin{multline*}
	 	\cX_{right} = \Big( \cX(\triv_{S_m};\phi_1)\times \cX(\triv_{S_m};\phi_2) \Big)\up_U^{T_w} \up_{T_w}^S \\
	 	= \Big( \cX(\triv_{S_m};\phi_1)\times \cX(\triv_{S_m};\phi_2) \Big)\up_U^{T_d} \up_{T_d}^S = \cX_{left},
	 \end{multline*}
	where the second equality follows from the transitivity of induction.
\end{proof}
